\documentclass[11pt,a4paper]{article}

\newcommand{\docTitle}{Plan Zarządzania Ryzykiem}

% ─────────────────────────────────────────────────────────────────
% Preambuła wspólna dla dokumentów projektu Feedback
% ─────────────────────────────────────────────────────────────────

% Kodowanie i język
\usepackage[utf8]{inputenc}
\usepackage[T1]{fontenc}
\usepackage[polish]{babel}

% Układ strony
\usepackage{geometry}
\geometry{margin=2.5cm}

% Tabele
\usepackage{booktabs}
\usepackage{tabularray}
\usepackage{longtable}
\usepackage{array}

% Grafika i strony specjalne
\usepackage{afterpage}
\usepackage{rotating}
\usepackage{pdflscape}

% Kolory
\usepackage{xcolor}
\definecolor{accent}{RGB}{30, 60, 100}
\definecolor{lowrisk}{RGB}{200, 230, 200}
\definecolor{medrisk}{RGB}{255, 235, 170}
\definecolor{highrisk}{RGB}{245, 195, 195}

% Tytuły, listy, akapity
\usepackage{titlesec}
\usepackage{parskip}
\usepackage{enumitem}
\usepackage{fancyhdr}

% Hyperref i cleveref – kolejność ma znaczenie:
% hyperref musi być załadowany przed cleveref
\usepackage{hyperref}
\hypersetup{colorlinks=true, linkcolor=accent, urlcolor=accent, citecolor=accent}
\usepackage[nameinlink]{cleveref}
% Polskie nazwy dla cleveref – formy skrócone unikają problemu z odmianą
\crefname{section}{sek.}{sek.}
\Crefname{section}{Sek.}{Sek.}
\crefname{subsection}{sek.}{sek.}
\Crefname{subsection}{Sek.}{Sek.}
\crefname{table}{tab.}{tab.}
\Crefname{table}{Tab.}{Tab.}
\crefname{figure}{rys.}{rys.}
\Crefname{figure}{Rys.}{Rys.}
\crefname{equation}{rów.}{rów.}
\Crefname{equation}{Rów.}{Rów.}
\crefname{enumi}{pkt}{pkt}
\Crefname{enumi}{Pkt}{Pkt}


% Numerowane sekcje z kolorem akcentu
\titleformat{\section}{\large\bfseries\color{accent}}{\thesection.}{0.6em}{}[\titlerule]
\titleformat{\subsection}{\normalsize\bfseries}{\thesubsection}{0.6em}{}
\titleformat{\subsubsection}{\normalsize\itshape}{\thesubsubsection}{0.6em}{}

% Łamanie wierszy – tolerancja dla polskich tekstów
\setlength{\emergencystretch}{3em}
\hyphenpenalty=1000
\tolerance=2000

% ─────────────────────────────────────────────────────────────────
% Nagłówki i stopki – wymaga zdefiniowania \docTitle przed \input
% ─────────────────────────────────────────────────────────────────

\pagestyle{fancy}
\fancyhf{}
\rhead{\small \docTitle}
\lhead{\small Projekt: \textit{Feedback}}
\cfoot{\thepage}
\renewcommand{\headrulewidth}{0.4pt}


\begin{document}
	\pagenumbering{Roman}
	% ── STRONA TYTUŁOWA ──────────────────────────────────────────────
	\begin{titlepage}
		\centering
		\vspace*{2cm}
		{\LARGE\bfseries\color{accent} \docTitle}\\[0.5cm]
		{\large Projekt: \textit{Feedback}}\\[0.3cm]
		{\normalsize Mechanika parowania jako klucz do sukcesu}\\[2cm]
		
		\begin{tabular}{ll}
			\textbf{Kierownik projektu:}   & Kacper Trzciński \\[0.3em]
			\textbf{Przewodnicząca KS:}    & Aleksandra Sobótka \\[0.3em]
			\textbf{Główny Dostawca:}      & Aleksander Płocharski \\[0.3em]
			\textbf{Główny Użytkownik:}    & Tomasz Herman \\[0.3em]
			\textbf{Zespół:}               & A. Bielicki, R. Głasek, J. Pączka, \\
			& M. Popkowicz, F. Shramko, K. Trzciński \\[0.3em]
			\textbf{Data opracowania:}     & \today \\[0.3em]
		\end{tabular}
		
		\vfill
		{\small Zarządzanie Przedsięwzięciami Informatycznymi\\
			Wydział Matematyki i Nauk Informacyjnych, Politechnika Warszawska}
	\end{titlepage}
	
	
	% ── SPIS TREŚCI ─────────────────────────────────────────────────
	\tableofcontents
	\pagenumbering{arabic}
	\newpage
	
	% ════════════════════════════════════════════════════════════════
	\section{Wstęp}\label{sec:wstep}
	% ════════════════════════════════════════════════════════════════
	
	\subsection{Przeznaczenie dokumentu}
	Niniejszy dokument opisuje sposób zarządzania ryzykiem w projekcie \textit{Feedback} ---
	grze \textit{puzzle-action platformer} tworzonej przez 6-osobowy zespół studencki w ramach
	przedmiotu \textit{Zarządzanie Przedsięwzięciami Informatycznymi}. Plan określa, jak ryzyko będzie
	oceniane, planowane i wdrażane, w jaki sposób będzie komunikowane oraz jakie tolerancje
	przyjmuje kierownictwo projektu.
	
	\subsection{Po co zarządzamy ryzykiem w projekcie?}
	Projekt prowadzony jest w dwóch semestrach przez zespół studencki obciążony równolegle
	innymi przedmiotami i pracą zawodową. Mechanika parowania, na której opiera się rozgrywka,
	wymaga precyzyjnej fizyki i wielu iteracji projektowych --- jej zbalansowanie nie jest
	gwarantowane. Bez systematycznego podejścia do ryzyka łatwo o utratę kontroli nad zakresem
	i terminem oddania. Zarządzanie ryzykiem ma na celu:
	\begin{itemize}[noitemsep]
		\item wczesne wykrywanie zagrożeń mogących wpłynąć na cele projektu;
		\item planowanie reakcji proporcjonalnych do ekspozycji projektu na dane ryzyko;
		\item utrzymanie ekspozycji projektu w ramach tolerancji zaakceptowanych przez Komitet Sterujący;
		\item wykorzystanie zidentyfikowanych szans (np. szybkiej iteracji prototypu) tam, gdzie zwiększają one wartość produktu.
	\end{itemize}
	
	\subsection{Definicje}
	\begin{description}[style=nextline, leftmargin=*]
		\item[Ryzyko] Niepewne zdarzenie lub zestaw zdarzeń, które --- jeśli wystąpią --- mogą mieć wpływ na osiągnięcie celów projektu. Ryzyko może być \textbf{zagrożeniem} (wpływ negatywny) albo \textbf{szansą} (wpływ pozytywny).
		\item[Prawdopodobieństwo] Szacowana częstotliwość lub szansa wystąpienia ryzyka w horyzoncie projektu.
		\item[Wpływ] Wielkość konsekwencji wystąpienia ryzyka dla celów projektu (zakres, harmonogram, jakość, ocena z przedmiotu).
		\item[Ekspozycja na ryzyko] Iloczyn prawdopodobieństwa i wpływu, wykorzystywany do priorytetyzacji ryzyk.
		\item[Tolerancja ryzyka] Akceptowalny poziom ekspozycji na ryzyko; przekroczenie tolerancji wymaga eskalacji do Komitetu Sterującego.
		\item[Właściciel ryzyka] Osoba odpowiedzialna za monitorowanie konkretnego ryzyka i wykonanie zaplanowanej reakcji.
	\end{description}
	
	% ════════════════════════════════════════════════════════════════
	\section{Ocena ryzyka}\label{sec:ocena}
	% ════════════════════════════════════════════════════════════════
	
	\subsection{Skala prawdopodobieństwa}
	Prawdopodobieństwo wystąpienia ryzyka oceniane jest w skali pięciostopniowej (patrz \cref{tab:prawdopodobienstwo}). Skala została zaprojektowana w sposób umożliwiający przypisanie wartości na podstawie obserwowalnych przesłanek w bieżącej sytuacji projektu, bez konieczności szacowania liczbowego prawdopodobieństwa.
	
	\begin{table}[ht!]
		\centering
		\begin{tblr}{
				colspec = {c l p{9cm}},
				hlines,
				row{1} = {font=\bfseries},
			}
			Ocena & Określenie & Charakterystyka \\
			1 & Bardzo niskie  & Wystąpienie ryzyka jest mało prawdopodobne; brak przesłanek w bieżącej sytuacji projektu. \\
			2 & Niskie         & Ryzyko może wystąpić w określonych okolicznościach; nie ma jeszcze obserwowanych sygnałów. \\
			3 & Średnie        & Wystąpienie ryzyka jest równie prawdopodobne co jego brak; występują pojedyncze przesłanki. \\
			4 & Wysokie        & Występują wyraźne przesłanki wskazujące na materializację ryzyka. \\
			5 & Bardzo wysokie & Materializacja ryzyka jest niemal pewna, jeżeli nie zostaną podjęte działania zaradcze. \\
		\end{tblr}
		\caption{Skala oceny prawdopodobieństwa wystąpienia ryzyka.}
		\label{tab:prawdopodobienstwo}
	\end{table}
	
	\subsection{Skala wpływu}
	Wpływ ryzyka oceniany jest w skali pięciostopniowej, z opisem konsekwencji dla kluczowych
	celów projektu: harmonogramu (terminowe oddanie dokumentów i wersji grywalnej), zakresu
	(kompletność wymagań funkcjonalnych) oraz oceny z przedmiotu (patrz
	\cref{tab:wplyw}).
	
	\begin{table}[ht!]
		\centering
		\begin{tblr}{
				colspec = {c l p{11cm}},
				hlines,
				row{1} = {font=\bfseries},
			}
			Ocena & Określenie & Charakterystyka wpływu \\
			1 & Bardzo niski & Wpływ pomijalny; opóźnienie do 2 dni roboczych, brak konieczności zmiany zakresu, brak wpływu na ocenę. \\
			2 & Niski        & Opóźnienie pojedynczego zadania o 3--5 dni roboczych, możliwe nadrobienie w bieżącym sprincie. \\
			3 & Średni       & Opóźnienie kamienia milowego o 1--2 tygodnie, drobna redukcja zakresu w obszarach \textit{Should} lub \textit{Could} (MoSCoW), pojedyncze uwagi prowadzącego. \\
			4 & Wysoki       & Opóźnienie kamienia milowego o ponad 2 tygodnie, konieczna redukcja zakresu obejmująca elementy \textit{Should}, istotne pogorszenie jakości oddania, ryzyko obniżenia oceny. \\
			5 & Bardzo wysoki & Zagrożenie terminu finalnej prezentacji, niezaliczenie etapu lub semestru, utrata kluczowej funkcjonalności (\textit{Must}). \\
		\end{tblr}
		\caption{Skala oceny wpływu ryzyka.}
		\label{tab:wplyw}
	\end{table}
	
	\subsection{Macierz ryzyka}\label{sec:macierz}
	Ekspozycja na ryzyko obliczana jest jako iloczyn prawdopodobieństwa i wpływu
	($E = P \cdot W$), co daje wartości w zakresie 1--25. Na podstawie ekspozycji ryzyka
	klasyfikuje się do trzech poziomów (patrz \cref{tab:macierz}):
	
	\begin{itemize}[noitemsep]
		\item \textbf{Niskie} ($1 \le E \le 5$, kolor zielony) --- akceptowane, monitorowane;
		\item \textbf{Średnie} ($6 \le E \le 12$, kolor żółty) --- wymagają zaplanowanej reakcji;
		\item \textbf{Wysokie} ($E \ge 15$, kolor czerwony) --- wymagają natychmiastowej reakcji i eskalacji do Komitetu Sterującego.
	\end{itemize}
	
	\begin{table}[ht!]
		\centering
		\begin{tblr}{
				colspec = {|c|c|c|c|c|c|},
				row{1} = {font=\bfseries, c},
				column{1} = {font=\bfseries},
				cells = {c},
				cell{2}{2}        = {bg=lowrisk},
				cell{2}{3}        = {bg=medrisk},
				cell{2}{4}        = {bg=highrisk},
				cell{2}{5}        = {bg=highrisk},
				cell{2}{6}        = {bg=highrisk},
				cell{3}{2}        = {bg=lowrisk},
				cell{3}{3}        = {bg=medrisk},
				cell{3}{4}        = {bg=medrisk},
				cell{3}{5}        = {bg=highrisk},
				cell{3}{6}        = {bg=highrisk},
				cell{4}{2}        = {bg=lowrisk},
				cell{4}{3}        = {bg=medrisk},
				cell{4}{4}        = {bg=medrisk},
				cell{4}{5}        = {bg=medrisk},
				cell{4}{6}        = {bg=highrisk},
				cell{5}{2}        = {bg=lowrisk},
				cell{5}{3}        = {bg=lowrisk},
				cell{5}{4}        = {bg=medrisk},
				cell{5}{5}        = {bg=medrisk},
				cell{5}{6}        = {bg=medrisk},
				cell{6}{2}        = {bg=lowrisk},
				cell{6}{3}        = {bg=lowrisk},
				cell{6}{4}        = {bg=lowrisk},
				cell{6}{5}        = {bg=lowrisk},
				cell{6}{6}        = {bg=lowrisk},
				hlines, vlines,
			}
			P \textbackslash{} W & 1 & 2 & 3 & 4 & 5 \\
			5 & 5  & 10 & 15 & 20 & 25 \\
			4 & 4  & 8  & 12 & 16 & 20 \\
			3 & 3  & 6  & 9  & 12 & 15 \\
			2 & 2  & 4  & 6  & 8  & 10 \\
			1 & 1  & 2  & 3  & 4  & 5  \\
		\end{tblr}
		
		\vspace{2mm}
		\caption{Macierz ekspozycji na ryzyko (prawdopodobieństwo $\times$ wpływ).}
		\label{tab:macierz}
	\end{table}
	
	\subsection{Kategorie reakcji na ryzyko}
	W projekcie stosowane są kategorie reakcji zgodne z metodyką PRINCE2, w podziale na
	zagrożenia i szanse (patrz \cref{tab:reakcje}).
	
	\begin{table}[ht!]
		\centering
		\begin{tblr}{
				colspec = {p{1.8cm} p{3cm} p{9.7cm}},
				hlines,
				row{1} = {font=\bfseries},
			}
			Typ ryzyka & Kategoria reakcji & Opis \\
			{\bfseries Zagrożenie} & Unikanie       & Eliminacja źródła ryzyka --- np. rezygnacja z elementu zakresu, który niesie nieproporcjonalne ryzyko. \\
			& Ograniczanie   & Redukcja prawdopodobieństwa lub wpływu --- np. wczesne prototypowanie, iteracyjne playtesty, redukcja zakresu w obszarach niekrytycznych. \\
			& Transfer       & Przeniesienie skutku finansowego lub merytorycznego na inną stronę. W projekcie studenckim zastosowanie ograniczone (brak dostawców komercyjnych). \\
			& Współdzielenie & Podział odpowiedzialności i konsekwencji z innym podmiotem (np. wspólne testowanie z zewnętrznymi playtesterami). \\
			& Akceptacja     & Świadome przyjęcie ryzyka bez aktywnej reakcji; możliwe utworzenie rezerwy czasowej. \\
			{\bfseries Szansa}     & Wykorzystanie  & Działania zwiększające prawdopodobieństwo do 100\% --- np. priorytetyzacja zadania, które otwiera dodatkową możliwość. \\
			& Zwiększenie    & Działania zwiększające prawdopodobieństwo lub wpływ szansy, ale nie dające pewności. \\
			& Współdzielenie & Włączenie partnera, który pomoże zrealizować szansę (np. zaproszenie testerów z innego rocznika do playtestów). \\
			& Akceptacja     & Przyjęcie szansy bez aktywnych działań --- skorzystanie z niej, jeśli wystąpi. \\
		\end{tblr}
		\caption{Kategorie reakcji na ryzyko stosowane w projekcie.}
		\label{tab:reakcje}
	\end{table}
	
	% ════════════════════════════════════════════════════════════════
	\section{Planowanie i wdrażanie reakcji na ryzyko}\label{sec:planowanie}
	% ════════════════════════════════════════════════════════════════
	
	\subsection{Wybór reakcji}
	Dobór reakcji odbywa się w oparciu o klasyfikację ekspozycji (\cref{sec:macierz}) oraz typ ryzyka
	(zagrożenie lub szansa). Ogólne zasady doboru reakcji:
	\begin{itemize}[noitemsep]
		\item Dla ryzyk \textbf{niskich} (zielonych) domyślną reakcją jest \textit{akceptacja} z monitorowaniem.
		\item Dla ryzyk \textbf{średnich} (żółtych) wymagana jest jedna z aktywnych reakcji (najczęściej \textit{ograniczanie}) z udokumentowanym planem działania.
		\item Dla ryzyk \textbf{wysokich} (czerwonych) wymagane są działania zaradcze równolegle z planem awaryjnym (\textit{fallback}); reakcja musi być uzgodniona z Komitetem Sterującym.
		\item Każde wybrane działanie musi mieć wskazanego \textbf{właściciela ryzyka}, termin wdrożenia oraz miernik skuteczności.
	\end{itemize}
	
	\subsection{Wdrażanie i weryfikacja skuteczności}
	Wdrożenie reakcji odbywa się w ramach bieżących sprintów na zasadach analogicznych do
	pozostałych zadań projektowych (rejestracja w systemie śledzenia zadań, oszacowanie czasu, przypisanie
	osoby). Po wdrożeniu działania właściciel ryzyka:
	\begin{enumerate}[noitemsep]
		\item ponownie ocenia prawdopodobieństwo i wpływ ryzyka po wdrożeniu reakcji (\textit{ryzyko rezydualne});
		\item aktualizuje pozycję w rejestrze ryzyk;
		\item zgłasza wynik kierownikowi projektu w trakcie najbliższego stand-upu.
	\end{enumerate}
	Jeżeli ryzyko rezydualne nadal przekracza tolerancję (\cref{sec:tolerancja}), planowane jest kolejne
	działanie zaradcze.
	
	\subsection{Plan awaryjny}
	Dla ryzyk o ekspozycji wysokiej i o krytycznym wpływie (5) opracowywany jest plan
	awaryjny --- zestaw działań uruchamianych w momencie materializacji ryzyka. Plan awaryjny
	zawiera: opis działań do podjęcia, osoby odpowiedzialne, szacowany czas wdrożenia oraz
	warunek aktywacji (zdarzenie, które uruchamia plan). Plany awaryjne są dokumentowane
	w rejestrze ryzyk jako oddzielne pole dla każdej pozycji o ekspozycji wysokiej.
	
	Harmonogram projektu nie przewiduje formalnej rezerwy czasowej --- termin finalnej
	prezentacji (początek semestru letniego 2026/2027) jest sztywno powiązany z kalendarzem
	akademickim. W praktyce oznacza to, że materializacja ryzyka wysokiego w późnej fazie
	projektu wymaga uruchomienia planu awaryjnego w trybie redukcji zakresu (zgodnie z
	priorytetyzacją MoSCoW), a nie wydłużenia harmonogramu.
	
	% ════════════════════════════════════════════════════════════════
	\section{Identyfikacja ryzyka}\label{sec:identyfikacja}
	% ════════════════════════════════════════════════════════════════
	
	\subsection{Źródła i odpowiedzialność}
	Za identyfikację ryzyka odpowiedzialny jest \textbf{cały zespół projektowy}.
	Każdy członek zespołu ma obowiązek zgłaszania zaobserwowanych zagrożeń i szans
	niezwłocznie po ich wykryciu. Kierownik projektu pełni rolę koordynatora --- konsoliduje
	zgłoszenia, weryfikuje ich kompletność i prowadzi rejestr ryzyk. Komitet Sterujący wnosi
	perspektywę zewnętrzną (biznesową i akademicką) i identyfikuje ryzyka niewidoczne z
	poziomu zespołu wykonawczego.
	
	\subsection{Techniki identyfikacji}
	W projekcie stosowane są następujące techniki identyfikacji ryzyk:
	\begin{description}[style=nextline]
		\item[Burza mózgów]
		Cykliczna burza mózgów dotycząca nowych ryzyk i szans odbywa się raz na sprint, podczas spotkania planującego kolejny sprint. Wszyscy członkowie zespołu zgłaszają obserwacje z minionego okresu.
		
		\item[SWOT]
		Analiza SWOT (Mocne strony / Słabe strony / Szanse / Zagrożenia) przeprowadzana jest na początku każdego z dwóch semestrów oraz przed kluczowymi kamieniami milowymi. Mocne i słabe strony zespołu wskazują obszary podatne na ryzyko wewnętrzne; szanse i zagrożenia dotyczą czynników zewnętrznych.
		
		\item[Wywiady]
		Krótkie indywidualne rozmowy kierownika projektu z członkami zespołu odbywają się przy każdej zmianie zakresu, składu lub harmonogramu. Wywiady ujawniają ryzyka, których członek zespołu mógłby nie zgłosić publicznie (przeciążenie, niejasność wymagań, problemy techniczne na początkowym etapie).
		
		\item[Przegląd dokumentacji]
		Analiza Karty projektu, Rejestru Interesariuszy, WBS, harmonogramu i wymagań pod kątem założeń i zależności, które mogą stać się ryzykiem.
		
		\item[Zgłoszenia ad hoc]
		Każdy członek zespołu może zgłosić ryzyko poza cyklem burzy mózgów --- przez kanał \texttt{\#ryzyka} na serwerze Discord zespołu lub bezpośrednią wiadomość do kierownika projektu.
	\end{description}
	
	\subsection{Narzędzia}
	\begin{itemize}[noitemsep]
		\item \textbf{Rejestr ryzyk} --- arkusz Google Sheets w folderze \texttt{/Dokumenty} na Google Drive; podstawowe narzędzie operacyjne.
		\item \textbf{Discord, kanał \texttt{\#ryzyka}} --- zgłaszanie ryzyk ad hoc oraz dyskusja nad reakcjami.
		\item \textbf{System śledzenia zadań w repozytorium kodu (z etykietą \texttt{risk})} --- techniczne ryzyka związane z konkretnym kodem lub komponentem; powiązane z odpowiednią pozycją w rejestrze.
	\end{itemize}
	
	\subsection{Utrzymanie aktualności rejestru ryzyk}
	Aby rejestr ryzyk nie utracił aktualności, w projekcie wprowadzono następujące mechanizmy:
	\begin{itemize}[noitemsep]
		\item \textbf{Przegląd cykliczny} --- pełny przegląd rejestru ryzyk odbywa się na końcu każdego sprintu, w ramach przeglądu sprintu z Komitetem Sterującym. Każda pozycja jest weryfikowana pod kątem aktualności prawdopodobieństwa, wpływu i statusu.
		\item \textbf{Aktualizacja ad hoc} --- każda zmiana stanu ryzyka (materializacja, zamknięcie, zmiana oceny) skutkuje natychmiastową aktualizacją rejestru przez właściciela ryzyka, najpóźniej w ciągu 48 godzin od wykrycia zmiany.
		\item \textbf{Audyt kierownika projektu} --- raz na miesiąc kierownik projektu przegląda rejestr pod kątem pozycji nieaktualizowanych dłużej niż 30 dni i konsultuje je z właścicielami.
	\end{itemize}
	
	\subsection{Struktura wpisu w rejestrze ryzyk}
	Każda pozycja w rejestrze ryzyk zawiera następujące pola:
	\begin{itemize}[noitemsep]
		\item identyfikator (R-XX);
		\item data zgłoszenia, autor zgłoszenia;
		\item typ (zagrożenie / szansa);
		\item opis zdarzenia, przyczyna, skutek;
		\item prawdopodobieństwo, wpływ, ekspozycja;
		\item kategoria reakcji i konkretne działanie zaradcze;
		\item właściciel ryzyka; termin wdrożenia działania;
		\item status (\textit{open / mitigated / materialised / closed});
		\item ryzyko rezydualne po wdrożeniu reakcji;
		\item historia zmian (data, zmiana, autor).
	\end{itemize}
	
	% ════════════════════════════════════════════════════════════════
	\section{Komunikacja zarządzania ryzykiem}\label{sec:komunikacja}
	% ════════════════════════════════════════════════════════════════
	
	Komunikacja w obszarze ryzyka stanowi część ogólnego systemu komunikacji opisanego w
	dokumencie \textit{Plan Zarządzania Komunikacją}. Poniżej wymieniono kanały i zdarzenia
	istotne specyficznie dla ryzyka.
	
	\subsection{Kanały i zdarzenia}
	\begin{description}[style=nextline]
		\item[Kanał \texttt{\#ryzyka} na Discord]
		Główny kanał do zgłaszania nowych ryzyk i dyskusji nad reakcjami. Czas reakcji kierownika projektu na zgłoszenie: do 48 godzin.
		
		\item[Stand-up tygodniowy]
		Punkt stały agendy: ,,nowe ryzyka i zmiany statusu''. Każdy członek zespołu zgłasza ewentualne nowe obserwacje lub zmiany w monitorowanych ryzykach.
		
		\item[Przegląd końca sprintu (z udziałem KS)]
		Pełen przegląd rejestru ryzyk; prezentacja statusu ryzyk wysokich, przegląd zmian w ocenie ryzyk średnich, omówienie ryzyk zmaterializowanych w sprincie. Decyzje dotyczące eskalacji.
		
		\item[Raport końca sprintu]
		Sekcja ,,Ryzyko'' --- lista nowych ryzyk, zmian w ocenie istniejących oraz ryzyk zmaterializowanych. Raport trafia do KS i prowadzących przez Discord serwera specjalności.
		
		\item[Eskalacja]
		Eskalacja do Komitetu Sterującego następuje niezwłocznie --- bez czekania na cykl raportowy --- w przypadku: (a) zidentyfikowania ryzyka o ekspozycji wysokiej (czerwonej); (b) materializacji ryzyka o wpływie 4 lub 5; (c) przekroczenia tolerancji ryzyka (\cref{sec:tolerancja}). Eskalacja prowadzona jest przez kierownika projektu w formie wiadomości na serwerze specjalności.
	\end{description}
	
	\subsection{Odpowiedzialności komunikacyjne}
	\begin{itemize}[noitemsep]
		\item \textbf{Każdy członek zespołu} --- zgłaszanie nowych ryzyk i zmian statusu ryzyk, w których jest właścicielem.
		\item \textbf{Kierownik projektu} --- agregacja zgłoszeń, aktualizacja rejestru, raportowanie do Komitetu Sterującego, eskalacja.
		\item \textbf{Komitet Sterujący} --- decyzje o akceptacji ryzyk wysokich, zatwierdzanie planów awaryjnych, decyzje w sytuacjach przekroczenia tolerancji.
	\end{itemize}
	
	% ════════════════════════════════════════════════════════════════
	\section{Tolerancja ryzyka}\label{sec:tolerancja}
	% ════════════════════════════════════════════════════════════════
	
	\subsection{Tolerancje na poziomie projektu}
	Poniższe tolerancje (\cref{tab:tolerancje}) stanowią propozycję zespołu projektowego sformułowaną w ramach
	niniejszego planu i podlegają zatwierdzeniu przez Komitet Sterujący w toku oceny
	dokumentu. Przekroczenie którejkolwiek z zaakceptowanych tolerancji w trakcie realizacji
	projektu skutkuje obowiązkiem eskalacji do KS i wymaga jego decyzji o dalszych krokach
	(dostosowanie zakresu, zmiana planu, akceptacja odstępstwa).
	
	\begin{table}[ht!]
		\centering
		\begin{tblr}{
				colspec = {p{4cm} p{10.5cm}},
				hlines,
				row{1} = {font=\bfseries},
			}
			Wymiar & Tolerancja \\
			Czas & Opóźnienie kamienia milowego o maksymalnie 1 tydzień bez eskalacji. Opóźnienie finalnej prezentacji nie jest dopuszczalne --- jest sztywno powiązane z początkiem semestru letniego 2026/2027. \\
			Zakres & Dopuszczalna jest redukcja elementów oznaczonych w priorytetyzacji MoSCoW jako \textit{Should} lub \textit{Could}. Redukcja elementów \textit{Must} (mechanika parowania, 3 grywalne poziomy, 2 typy przeciwników, prezentacja przed nowym rocznikiem) wymaga decyzji KS. \\
			Jakość & Dopuszczalne jest odstępstwo od pojedynczych wymagań niefunkcjonalnych w wymiarze do 20\% (np. spadek minimalnych FPS z 60 do 48). Większe odstępstwa wymagają decyzji KS. \\
			Budżet & Budżet projektu wynosi 0~zł i nie podlega zmianom. Każda potencjalna potrzeba poniesienia kosztu finansowego wymaga eskalacji. \\
			Ryzyko (zagregowane) & Liczba aktywnych ryzyk wysokich (czerwonych) w rejestrze nie powinna przekraczać 3 jednocześnie. Przekroczenie tej liczby wymaga przeglądu nadzwyczajnego z KS. \\
		\end{tblr}
		\caption{Tolerancje na poziomie projektu.}
		\label{tab:tolerancje}
	\end{table}
	
	\subsection{Tolerancje pojedynczego ryzyka}
	\begin{itemize}[noitemsep]
		\item \textbf{Ryzyko niskie} ($1 \le E \le 5$) --- akceptowane bez ograniczeń; monitorowanie cykliczne.
		\item \textbf{Ryzyko średnie} ($6 \le E \le 12$) --- akceptowane pod warunkiem wdrożenia zaplanowanej reakcji; brak wdrożenia w terminie powoduje eskalację.
		\item \textbf{Ryzyko wysokie} ($E \ge 15$) --- niemożliwe do zaakceptowania bez decyzji KS; wymaga planu awaryjnego.
	\end{itemize}
	
	\subsection{Postępowanie w przypadku przekroczenia tolerancji}
	W przypadku przekroczenia tolerancji kierownik projektu:
	\begin{enumerate}[noitemsep]
		\item niezwłocznie informuje Przewodniczącą Komitetu Sterującego (kanał: Discord serwera specjalności);
		\item w ciągu 48 godzin przygotowuje krótki dokument zawierający opis sytuacji, możliwe scenariusze działania i rekomendację;
		\item Komitet Sterujący podejmuje decyzję o jednym z scenariuszy: redukcja zakresu, modyfikacja harmonogramu, akceptacja odstępstwa lub zamknięcie etapu;
		\item decyzja KS jest dokumentowana w rejestrze ryzyk i komunikowana zespołowi w ciągu 24 godzin.
	\end{enumerate}
	
	% ════════════════════════════════════════════════════════════════
	\section{Przegląd planu}\label{sec:przeglad}
	% ════════════════════════════════════════════════════════════════
	
	\noindent Niniejszy plan podlega przeglądowi na końcu każdego etapu projektu oraz przy każdej
	istotnej zmianie kontekstu (zmiana składu zespołu, zmiana zakresu, zmiana prowadzącego
	przedmiot). W razie potrzeby modyfikacji kierownik projektu aktualizuje dokument, opisuje
	zmianę w sekcji historii i informuje wszystkich interesariuszy zgodnie z planem komunikacji.
	
\end{document}
